%&template
\documentclass{article}
\usepackage{src/template}

\endofdump
\usepackage{minted}

\date{2026 -- 01 -- 11}
\useTemplate[NPA, Exercise=10]
\title{NPA -- Exercise 10}


\begin{document}
\maketitle
    
    \section{}
    \subsection{}

        The motivation behind \texttt{Scenic routing} is to route packets that are normally routed through dark fiber or narrow copper lines to their destination via more appropriate transport routes such as air freight or radio. A higher purpose would be to route them preferentially along culturally important locations around the world in order to promote intercultural acceptance and understanding of other cultures.
        
        However, the main goal is to get more fresh air and sunshine for the packets.
    
    \subsection{}

        \texttt{Scenic routing} should be enabled using the IPv6 hop-by-hop option header. This option header is used to give a network device extended instructions on how to process the packet. The identifier $0x0A$ should be used here. Network nodes must process such marked packets according to scenic routing. The \texttt{Scenic routing Options} (SRO) use 8-bit coding to specify whether the packet must or can be routed via airborne paths. It also specifies whether this packet is afraid of flying or can tolerate radiation, i.e., which routes are possible. A preference can also be specified. Network providers must provide or supplement this infrastructure as far as possible.

        In order to transfer a package by air, the network devices must perform many additional functions. When transported as air freight, e.g., by pigeon, the packets have to wait longer before boarding the animal carrier. This requires a queue. Since the destination is fresh air and sunshine, the queue should perhaps be equipped with a sun terrace. Even though more sunshine is desirable for our hard-working packets, we also have to be realistic. The packages work in IT and rarely tolerate a lot of sunscreen. Skin cancer is a serious issue, which is why network devices should regularly apply sunscreen to the bits.
    
    \subsection{}

        We find this approach very innovative, and improving working conditions for packages is a worthwhile goal. More relaxed bits lead to better morale and thus possibly fewer bit flips. However, if the bits suffer sunburn, this does not help the healthcare system either; skin cancer is our biggest concern. If the bits fail due to health reasons, a recovery phase is added to the longer transport time from the sender to the destination. This is not necessarily something that customers would appreciate. Basically, this is a matter of weighing up the pros and cons and must be decided on a case-by-case basis and also in conjunction with the packages.
    
    


\end{document}